\documentclass[11pt,letterpaper]{article}
\usepackage[margin=0.65in]{geometry}
\usepackage{fontspec}
\setmainfont{Times New Roman}
\setmonofont{Consolas}[Scale=0.85]
\usepackage{booktabs}
\usepackage[hidelinks]{hyperref}
\pagestyle{empty}
\setlength{\parindent}{0pt}
\setlength{\parskip}{5pt}
\setlength{\emergencystretch}{1em}
\newcommand{\topic}[1]{\par\smallskip\textbf{#1}\par\nopagebreak}
\begin{document}
\begin{center}
{\Large\bfseries Lab 3 Report}\\[4pt]
CSE 60685 Machine Learning for Embedded Systems\\[2pt]
Structured Pruning and Fine-Tuning on Raspberry Pi
\end{center}

\textbf{Device and source.} Supplied Raspberry Pi 4B Rev.\ 1.5, Raspberry Pi OS 64-bit (Debian 13), Python 3.13.5 and PyTorch 2.8.0+cpu. All experiments ran on the Pi CPU, using the CanaKit supply, closed case, running fan and \texttt{ondemand} governor. I used my own final five-epoch \textbf{Lab 2 A} checkpoint (44,426 parameters), not the fallback. Import verification reproduced its 79.65\% validation accuracy. All instructor code was unchanged.

\textbf{Common settings.} The supplied Fashion-MNIST split contains 12,000 training and 2,000 validation images; final evaluation uses all 10,000 test images. Inputs are $28\times28$ grayscale float32 in $[0,1]$, without augmentation or extra normalization. All branches start independently from the same imported source. Each, including Control, receives \textbf{three additional epochs}, batch 64, a new Adam optimizer at 0.0001, seed 42, one intra-op and one inter-op thread, and zero loader workers. The last-epoch checkpoints were reloaded and verified before final test evaluation.

\begin{center}
\begin{tabular}{lrrrr}
\toprule
Model & Parameters & \shortstack{Validation before $\rightarrow$ after\\fine-tuning (\%)} & \shortstack{Test\\(\%)} & \shortstack{Median inference\\(ms)}\\
\midrule
Control & 44,426 & 79.65 $\rightarrow$ 82.60 & 81.77 & 2.6906 \\
S25 & 36,142 & 77.45 $\rightarrow$ 82.35 & 81.70 & 2.6028 \\
S50 & 27,858 & 64.05 $\rightarrow$ 81.40 & 81.00 & 2.5076 \\
\bottomrule
\end{tabular}
\end{center}

\textbf{Timing.} All training and test evaluation finished before sequential benchmarks. Each used the same preloaded validation image, shape $[1,1,28,28]$, evaluation/inference mode, one CPU thread, 10 untimed warmups and 100 recorded calls. The unchanged timer covers synchronous \texttt{model(tensor)} plus Python call overhead; loading, preprocessing, argmax/softmax and file output are excluded. Temperatures before warmup were 41.4--43.3\,$^{\circ}$C; all sampled throttling flags were \texttt{0x0}.

\topic{1. Pruning, parameter count and accuracy recovery}
S25/S50 remove 25\%/50\% of \emph{Conv2 output channels}, leaving 12/8 of 16. The supplied L1 ranking retains the largest filter scores, copies their filters and biases, and selects the matching blocks of 16 FC1 input columns. FC1 input widths become 192/128 instead of 256; Conv1 and hidden widths 120/84 stay unchanged. Both compact-model checks passed against the matching masked-feature reference.

Relative to the continued-training Control, S25 and S50 reduce total parameters by \textbf{18.65\%} and \textbf{37.29\%}. Their validation accuracies recover by \textbf{4.90} and \textbf{17.35 percentage points} from immediately after pruning. Control improves by 2.95 points with the same extra training. Final validation is 0.25/1.20 points below Control for S25/S50; final test accuracy is only 0.07/0.77 points below Control. Thus, fine-tuning substantially repairs pruning damage, while the larger cut retains a larger accuracy gap.

\topic{2. Measured speed and deployment choice}
Parameter reduction does not translate proportionally into latency reduction: S25 and S50 reduce the median by \textbf{3.26\% and 6.80\%}, respectively, yielding 1.03$\times$ and 1.07$\times$ speedups (Control median divided by candidate median). Inference covers the whole network. Conv1, later classifier layers and execution overhead remain; parameters also differ in how often they are used across spatial positions. CPU kernel efficiency and memory access can further affect runtime. These aggregate measurements do not identify an individual layer as the bottleneck.

I choose \textbf{S50} because minimizing parameter count is my priority. It has the fewest parameters (27,858) and the lowest measured median (2.5076 ms), with test accuracy of 81.00\%. I prefer this smaller model while accepting the measured 0.77-percentage-point accuracy reduction relative to Control. These are single-run comparisons with 100 inference calls per model; the accuracy gaps are not a statistical equivalence claim.

\vfill
{\footnotesize Sources: \href{https://canvas.nd.edu/courses/144914/files/6848884?wrap=1}{Lab 3 Tutorial}; \href{https://canvas.nd.edu/courses/144914/files/6848882?wrap=1}{Assignment}; \href{https://github.com/guoyb17/CSE60685-FA26-Lab-3/tree/4c635083c39862381940db06d134bfe95ce92088}{instructor repository, \texttt{4c63508}}. Original checkpoints, raw JSON/CSV, source hashes and execution logs are retained.}

\newpage
\begin{center}
{\Large\bfseries AI Attribution Appendix}\\[4pt]
Lab 3
\end{center}
\textbf{Tool.} OpenAI Codex.

\textbf{Main requests} (summarized): help complete the lab according to its Tutorial and Assignment, troubleshoot the Pi connection, and prepare a concise LaTeX report with retained evidence.

\textbf{Assistance.} Codex helped interpret the instructions, run the supplied commands through SSH, and create separate helpers for execution, cooling, logging and evidence checks. It assisted with numerical comparisons, explanation, English wording and LaTeX formatting. All instructor model, pruning, training and timing code remained unchanged.

\textbf{Student contribution and evidence.} I set up the hardware and provided my deployment preference and reasoning. All reported measurements came from the actual Raspberry Pi; original checkpoints, JSON/CSV, logs and the detailed assistance record are retained.
\end{document}
